% Unit 8 review sheet -- 2 pages.
\documentclass{apchem-worksheet}

\apcunit{8}
\apcunittitle{Acids and Bases}
\apcdoctitle{Review Sheet}
\apcsubtitle{Everything on the exam traces to this page \textbullet\ exam Fri Feb 19}

\begin{document}
\apctitleblock

\section*{\sffamily First: which of the four cases are you in?}

\renewcommand{\arraystretch}{1.3}
\noindent\begin{tabular}{@{}p{4.4cm}p{9.3cm}@{}}
\toprule
\textbf{Strong acid or base alone} & complete ionization; pH directly.
  Double $[\ce{OH-}]$ for \ce{Ba(OH)2}, \ce{Ca(OH)2}, \ce{Sr(OH)2}\\
\textbf{Weak acid or base alone} & one equilibrium.
  $[\ce{H+}] = \sqrt{K_aC_0}$, or for a base
  $[\ce{OH-}] = \sqrt{K_bC_0}$ --- then convert through pOH\\
\textbf{Both members of a pair} & buffer.
  pH $=$ p$K_a + \log(\text{base}/\text{acid})$\\
\textbf{Strong reagent added} & react to completion \emph{first}, in moles;
  then whichever case is left\\
\bottomrule
\end{tabular}

\section*{\sffamily The numbers and relationships}

\begin{itemize}[leftmargin=*,itemsep=0.16em]
  \item $\mathrm{pH} + \mathrm{pOH} = 14.00$;
        $K_aK_b = K_w = 1.0\times10^{-14}$;
        p$K_a = -\log K_a$.
  \item Six strong acids: \ce{HCl}, \ce{HBr}, \ce{HI}, \ce{HNO3},
        \ce{H2SO4}, \ce{HClO4}. \ce{HF} is \emph{not} one.
  \item Neutral means $[\ce{H+}] = [\ce{OH-}]$ --- \textbf{not} pH 7.
  \item Optimal buffering at pH $=$ p$K_a$; useful range p$K_a \pm 1$;
        indicators change over p$K_a \pm 1$ too.
\end{itemize}

\section*{\sffamily Titration curve, region by region}

\begin{center}
\renewcommand{\arraystretch}{1.15}
\begin{tabular}{@{}lll@{}}
\toprule
\textbf{Region} & \textbf{Major species} & \textbf{Method} \\
\midrule
before any base & \ce{HA} & weak-acid ICE \\
before equivalence & \ce{HA} $+$ \ce{A-} & buffer, H--H \\
\textbf{halfway} & equal amounts & \textbf{pH $=$ p$K_a$} \\
at equivalence & \ce{A-} only & weak-base ICE, $K_b = K_w/K_a$ \\
after equivalence & excess \ce{OH-} & strong base alone \\
\bottomrule
\end{tabular}
\end{center}

Equivalence pH: strong$+$strong $= 7.00$ exactly; weak acid $+$ strong base
$>7$; weak base $+$ strong acid $<7$.

\section*{\sffamily The traps, one line each}

\begin{itemize}[leftmargin=*,itemsep=0.2em]
  \item Reporting $-\log x$ as the pH for a weak \emph{base} --- it is the
        pOH.
  \item Forgetting to double $[\ce{OH-}]$ for a Group 2 hydroxide.
  \item Skipping the 5\% check, or running it and ignoring the result.
  \item Putting added strong base into an ICE table instead of reacting it
        to completion first.
  \item Inverting the base/acid ratio in Henderson--Hasselbalch.
  \item Assuming any equivalence point is at pH 7.
  \item \emph{Calculating} solubility from pH --- topic 8.11 excludes it.
\end{itemize}

\section*{\sffamily Self-check --- 12 questions, exam conditions}

\begin{enumerate}[leftmargin=*,itemsep=0.34em]
  \item pH of \SI{0.010}{\Molar} \ce{Ba(OH)2}:
        \answerblank[2cm]{12.30}
  \item pH of \SI{0.10}{\Molar} acetic acid ($K_a = 1.8\times10^{-5}$):
        \answerblank[2cm]{2.87}
  \item pH of \SI{0.100}{\Molar} \ce{NH3} ($K_b = 1.8\times10^{-5}$):
        \answerblank[2cm]{11.13}
  \item $K_b$ for acetate: \answerblank[3.2cm]{$5.6\times10^{-10}$}
  \item \ce{NH4Cl} solutions are: \answerblank[2cm]{acidic}
  \item pH of a buffer with equal amounts of acetic acid and acetate:
        \answerblank[2cm]{4.74}
  \item Buffer at pH 5.22 from acetic acid: base/acid ratio?
        \answerblank[2cm]{3.0}
  \item pH at the halfway point of a weak acid titration:
        \answerblank[2cm]{p$K_a$}
  \item Equivalence pH for acetic acid $+$ \ce{NaOH}
        (\SI{0.10}{\Molar} each, \SI{50}{\milli\litre}):
        \answerblank[2cm]{8.72}
  \item Equivalence pH for \ce{HCl} $+$ \ce{NaOH}:
        \answerblank[2cm]{7.00}
  \item Indicator for an equivalence pH of 8.7:
        \answerblank[3.4cm]{phenolphthalein}
  \item Is \ce{AgCl} more soluble in acid?
        \answerblank[1.6cm]{no}
\end{enumerate}

\begin{apcnote}[How to study for this exam]
This is a 44-point exam covering nine blocks --- the second-heaviest unit in
the course. Budget more time than for any unit except 3.

Redo WS 5's FRQ cold; the long free-response is its twin, and it
deliberately walks the whole titration curve. Then drill the \emph{method
selection} table on page one until naming the case is automatic, because
almost every lost mark in this unit comes from running the right arithmetic
on the wrong situation.

Two sanity checks worth thirty seconds each: for a base problem, does your
pH come out above 7? And after a buffer calculation, is the pH within one
unit of p$K_a$? If either fails, you have almost certainly picked the wrong
case.
\end{apcnote}

\end{document}
